\documentclass[11pt]{article}
\usepackage[utf8]{inputenc}
\usepackage{amsmath, amssymb, amsthm}
\usepackage{geometry}
\geometry{a4paper, margin=1in}
\usepackage{enumitem}

\title{\textbf{Homework 16.4: Green's Theorem\\Problem Set \& Solutions}}
\author{Comprehensive Diagnostic and Practice}
\date{}

\begin{document}

\maketitle

\tableofcontents
\vspace{1cm}

\section{Part 1: Problem Set}

\textbf{Directions:} For each of the following problems, evaluate the line integral or physical quantity using Green's Theorem. Pay close attention to orientations, singularities, and opportunities for algebraic shortcuts. 

\begin{enumerate}[label=\textbf{Q\arabic*.}, leftmargin=*]
    % Concept 1: Verification
    \item Verify Green's Theorem for the line integral $\oint_C (3x - y) \,dx + (x + 5y) \,dy$ where $C$ is the boundary of the unit square $[0,1] \times [0,1]$ oriented counterclockwise. Evaluate both the line integral directly and the double integral.
    \item Verify Green's Theorem for $\oint_C xy^2 \,dx + x^2 y \,dy$ where $C$ is the counterclockwise triangle with vertices $(0,0), (2,0), (0,2)$.
    
    % Concept 2: Standard Cartesian
    \item Evaluate $\oint_C (y^2 + \sin(x^2)) \,dx + (x^2 + \cos(y^2)) \,dy$ where $C$ is the boundary of the region enclosed by $y = x$ and $y = x^2$, oriented counterclockwise.
    \item Evaluate $\oint_C (e^x - 3y^2) \,dx + (e^y + 4xy) \,dy$ where $C$ is the counterclockwise boundary of the rectangle $(0,0), (4,0), (4,2), (0,2)$.
    \item Evaluate $\oint_C (x^4 + y) \,dx + (2xy) \,dy$ where $C$ is the counterclockwise boundary of the region bounded by $y = x^2$ and $y = 4$.
    \item Evaluate $\oint_C y^3 \,dx - x^3 \,dy$ where $C$ is the square $[0,2] \times [0,2]$ oriented counterclockwise.
    \item Evaluate $\oint_C (x \sin x - y^2) \,dx + (x + y \cos y) \,dy$ along the triangle with vertices $(0,0), (1,1), (0,1)$ oriented counterclockwise.
    
    % Concept 3: Polar Transformation
    \item Evaluate $\oint_C -y^3 \,dx + x^3 \,dy$ over the upper half of the circle $x^2 + y^2 = 4$ and the line segment from $(-2,0)$ to $(2,0)$, oriented counterclockwise.
    \item Evaluate $\oint_C (5y + \tan^{-1}x) \,dx + (2x + e^{y^2}) \,dy$ over the positively oriented annular region bounded by $x^2 + y^2 = 1$ and $x^2 + y^2 = 4$.
    \item Evaluate $\oint_C (3x - 5y) \,dx + (x + 4y) \,dy$ over the circle $x^2 + y^2 = 16$ oriented counterclockwise.
    \item Evaluate $\oint_C -y \,dx + x \,dy$ over the counterclockwise boundary of the quarter circle $x^2 + y^2 \le 25$ in the first quadrant.
    \item Evaluate $\oint_C (e^x - y^3) \,dx + (\cos y + x^3) \,dy$ where $C$ is the boundary of the semi-circle $x^2 + y^2 \le 9$ for $y \ge 0$, oriented counterclockwise.
    
    % Concept 4: Negative/Clockwise Orientation
    \item Evaluate $\oint_C y^2 \,dx + x^2 \,dy$ where $C$ is the rectangle traversed in the order $(0,0) \to (0,3) \to (2,3) \to (2,0) \to (0,0)$.
    \item Evaluate $\oint_C y \,dx - x \,dy$ where $C$ is the unit circle $x^2 + y^2 = 1$ traversed clockwise.
    \item Evaluate $\oint_C (\sin x + y) \,dx + (x^2 + \cos y) \,dy$ around the triangle path $(0,0) \to (0,1) \to (1,0) \to (0,0)$.
    \item Evaluate $\oint_C (x^3 - y) \,dx + (x + y^3) \,dy$ around the circle $x^2 + y^2 = 4$ oriented clockwise.
    
    % Concept 5: Area Shortcut
    \item Let $C$ be any closed, simply-connected curve oriented counterclockwise that encloses an area of 10. Evaluate $\oint_C (7y - e^{\sin x}) \,dx + (15x - \ln(y^2+1)) \,dy$.
    \item Let $D$ be a region with an area of 5, and $C$ be its positively oriented boundary. Evaluate $\oint_C (3y^2 - 4y) \,dx + (2x + 6xy) \,dy$.
    \item Find the work done by the force field $\mathbf{F} = \langle \ln(1+x^2) - 2y, 5x + \sin(\sqrt{y}) \rangle$ on a particle moving counterclockwise once around the unit circle.
    
    % Concept 6: Geometric Properties
    \item Use Green's Theorem formula $A = \frac{1}{2} \oint_C (x \,dy - y \,dx)$ to find the area of the polygon with vertices $(0,0), (2,0), (3,3), (1,4)$.
    \item Find the centroid coordinate $\bar{x}$ of the upper half of the unit circle $x^2+y^2 \le 1, y \ge 0$ using Green's Theorem identity $\bar{x} = \frac{1}{2A} \oint_C x^2 \,dy$.
    \item Use Green's Theorem to calculate the area bounded by the ellipse $\frac{x^2}{9} + \frac{y^2}{16} = 1$.
    \item Use the moment of inertia formula $I_x = -\frac{\rho}{3} \oint_C y^3 \,dx$ to find $I_x$ for the uniform rectangular region $[0, a] \times [0, b]$ with density $\rho=1$.
    
    % Concept 7: Extended Theorem / Singularities
    \item Evaluate $\oint_C \frac{-y}{x^2+y^2} \,dx + \frac{x}{x^2+y^2} \,dy$ where $C$ is the square with vertices $(\pm 2, \pm 2)$ oriented counterclockwise.
    \item Evaluate $\oint_C \left(\frac{-y}{x^2+y^2} + 2y\right) \,dx + \left(\frac{x}{x^2+y^2} + 3x\right) \,dy$ over the circle $x^2+y^2 = 9$ oriented counterclockwise.
    \item Evaluate $\oint_C \mathbf{F} \cdot d\mathbf{r}$ where $\mathbf{F} = \left\langle \frac{-y}{(x-1)^2+y^2}, \frac{x-1}{(x-1)^2+y^2} \right\rangle$ and $C$ is the circle $x^2+y^2=16$ oriented counterclockwise.
    
    % Concept 8: Odd Functions / Symmetries
    \item Evaluate $\oint_C (\sin(x^3) - y) \,dx + (x^2 y^2 + x) \,dy$ where $C$ is the boundary of the rectangle $[-1,1] \times [0,2]$ oriented counterclockwise.
    \item Evaluate $\oint_C (x^4 + y^2) \,dx + (x^3 y + y^4) \,dy$ over the region bounded by the parabolas $x = -y^2$ and $x = y^2 - 8$, oriented counterclockwise.
    
    % Concept 9: Untested / Advanced Parameterization
    \item Find the area of the region bounded by the astroid $x = a \cos^3 t$, $y = a \sin^3 t$, for $0 \le t \le 2\pi$ using Green's Theorem area formulas.
    \item Find the area under one arch of the cycloid given by $x = R(t - \sin t)$, $y = R(1 - \cos t)$ for $t \in [0, 2\pi]$ and above the x-axis, using Green's Theorem.
\end{enumerate}

\newpage
\section{Part 2: Detailed Solutions}

\begin{itemize}[leftmargin=*]
    \item \textbf{Solution 1:} $P = 3x - y$, $Q = x + 5y \implies Q_x - P_y = 1 - (-1) = 2$.
    \textit{Double Integral:} $\iint_D 2 \,dA = 2 \times \text{Area} = 2(1) = 2$.
    \textit{Line Integral:} 
    Bottom (y=0, dy=0): $\int_0^1 3x \,dx = 3/2$.
    Right (x=1, dx=0): $\int_0^1 (1 + 5y) \,dy = 1 + 5/2 = 7/2$.
    Top (y=1, dy=0, x from 1 to 0): $\int_1^0 (3x - 1) \,dx = -[3/2 - 1] = -1/2$.
    Left (x=0, dx=0, y from 1 to 0): $\int_1^0 5y \,dy = -5/2$.
    Sum $= 3/2 + 7/2 - 1/2 - 5/2 = 2$. Both evaluate to \textbf{2}.

    \item \textbf{Solution 2:} $P = xy^2$, $Q = x^2 y \implies Q_x - P_y = 2xy - 2xy = 0$. By Green's Theorem, $\iint_D 0 \,dA = \textbf{0}$.

    \item \textbf{Solution 3:} $P = y^2 + \sin(x^2)$, $Q = x^2 + \cos(y^2) \implies Q_x - P_y = 2x - 2y$.
    The region $D$ is $0 \le x \le 1$, $x^2 \le y \le x$.
    $\int_0^1 \int_{x^2}^x (2x - 2y) \,dy \,dx = \int_0^1 [2xy - y^2]_{x^2}^x \,dx = \int_0^1 ((2x^2 - x^2) - (2x^3 - x^4)) \,dx = \int_0^1 (x^2 - 2x^3 + x^4) \,dx = \frac{1}{3} - \frac{2}{4} + \frac{1}{5} = \frac{10 - 15 + 6}{30} = \mathbf{\frac{1}{30}}$.

    \item \textbf{Solution 4:} $P = e^x - 3y^2$, $Q = e^y + 4xy \implies Q_x - P_y = 4y - (-6y) = 10y$.
    $\int_0^4 \int_0^2 10y \,dy \,dx = \int_0^4 [5y^2]_0^2 \,dx = \int_0^4 20 \,dx = \textbf{80}$.

    \item \textbf{Solution 5:} $Q_x - P_y = 2y - 1$. Region $D$: $-2 \le x \le 2$, $x^2 \le y \le 4$.
    $\int_{-2}^2 \int_{x^2}^4 (2y - 1) \,dy \,dx = \int_{-2}^2 [y^2 - y]_{x^2}^4 \,dx = \int_{-2}^2 (12 - x^4 + x^2) \,dx$.
    Since the integrand is even, $2 \int_0^2 (12 + x^2 - x^4) \,dx = 2 [12x + \frac{x^3}{3} - \frac{x^5}{5}]_0^2 = 2(24 + \frac{8}{3} - \frac{32}{5}) = 2(\frac{360 + 40 - 96}{15}) = \mathbf{\frac{608}{15}}$.

    \item \textbf{Solution 6:} $Q_x - P_y = -3x^2 - 3y^2 = -3(x^2 + y^2)$.
    $\int_0^2 \int_0^2 -3(x^2 + y^2) \,dy \,dx = -3 \int_0^2 [x^2y + \frac{y^3}{3}]_0^2 \,dx = -3 \int_0^2 (2x^2 + \frac{8}{3}) \,dx = -3 [\frac{2x^3}{3} + \frac{8x}{3}]_0^2 = -3 (\frac{16}{3} + \frac{16}{3}) = \textbf{-32}$.

    \item \textbf{Solution 7:} $Q_x - P_y = 1 - (-2y) = 1 + 2y$. Region $D$: $0 \le x \le 1$, $x \le y \le 1$.
    $\int_0^1 \int_x^1 (1 + 2y) \,dy \,dx = \int_0^1 [y + y^2]_x^1 \,dx = \int_0^1 (2 - x - x^2) \,dx = 2 - \frac{1}{2} - \frac{1}{3} = \mathbf{\frac{7}{6}}$.

    \item \textbf{Solution 8:} $Q_x - P_y = 3x^2 - (-3y^2) = 3(x^2 + y^2) = 3r^2$.
    Using polar coordinates for the half-circle (radius 2): $\int_0^\pi \int_0^2 3r^2(r) \,dr \,d\theta = 3 \int_0^\pi d\theta \int_0^2 r^3 \,dr = 3(\pi)(4) = \mathbf{12\pi}$.

    \item \textbf{Solution 9:} $Q_x - P_y = 2 - 5 = -3$. 
    The double integral is $-3 \iint_D dA = -3 \times (\text{Area of Annulus}) = -3(\pi(2^2) - \pi(1^2)) = -3(3\pi) = \mathbf{-9\pi}$.

    \item \textbf{Solution 10:} $Q_x - P_y = 1 - (-5) = 6$. 
    $6 \iint dA = 6 \times (\text{Area of circle}) = 6(16\pi) = \mathbf{96\pi}$.

    \item \textbf{Solution 11:} $Q_x - P_y = 1 - (-1) = 2$.
    $2 \iint dA = 2 \times (\text{Area of Quarter Circle}) = 2(\frac{25\pi}{4}) = \mathbf{\frac{25\pi}{2}}$.

    \item \textbf{Solution 12:} $Q_x - P_y = 3x^2 - (-3y^2) = 3r^2$.
    Semi-circle radius 3: $\int_0^\pi \int_0^3 3r^2(r) \,dr \,d\theta = \pi [3 \frac{r^4}{4}]_0^3 = \mathbf{\frac{243\pi}{4}}$.

    \item \textbf{Solution 13:} \textit{Clockwise Orientation Trap.} Multiply integral by $-1$.
    $Q_x - P_y = 2x - 2y$. $-1 \times \int_0^2 \int_0^3 (2x - 2y) \,dy \,dx = \int_0^2 \int_0^3 (2y - 2x) \,dy \,dx = \int_0^2 [y^2 - 2xy]_0^3 \,dx = \int_0^2 (9 - 6x) \,dx = [9x - 3x^2]_0^2 = 18 - 12 = \textbf{6}$.

    \item \textbf{Solution 14:} \textit{Clockwise.} $Q_x - P_y = -1 - 1 = -2$.
    Standard integral evaluates to $-2(\text{Area}) = -2\pi$. Since orientation is clockwise, multiply by $-1 \implies \mathbf{2\pi}$.

    \item \textbf{Solution 15:} \textit{Clockwise trace.} Path traces $(0,0)$ up to $(0,1)$, down-right to $(1,0)$, left to $(0,0)$. This is clockwise.
    $Q_x - P_y = 2x - 1$. Apply $-1$: $\int_0^1 \int_0^{1-x} (1 - 2x) \,dy \,dx = \int_0^1 (1 - 2x)(1 - x) \,dx = \int_0^1 (1 - 3x + 2x^2) \,dx = [x - \frac{3x^2}{2} + \frac{2x^3}{3}]_0^1 = 1 - 1.5 + 0.666... = \mathbf{\frac{1}{6}}$.

    \item \textbf{Solution 16:} \textit{Clockwise.} $Q_x - P_y = 1 - (-1) = 2$.
    $2 \times (4\pi) = 8\pi$. Multiply by $-1 \implies \mathbf{-8\pi}$.

    \item \textbf{Solution 17:} Area shortcut. $Q_x - P_y = 15 - 7 = 8$. 
    Answer $= 8 \times \text{Area} = 8(10) = \textbf{80}$.

    \item \textbf{Solution 18:} $Q_x = 2 + 6y$, $P_y = 6y - 4$. $Q_x - P_y = (2+6y) - (6y-4) = 6$.
    Integral $= 6 \times \text{Area} = 6(5) = \textbf{30}$.

    \item \textbf{Solution 19:} Work $= \oint \mathbf{F} \cdot d\mathbf{r}$. $Q_x = 5$, $P_y = -2$. $Q_x - P_y = 7$.
    Area of unit circle is $\pi$. Work $= 7 \times \pi = \mathbf{7\pi}$.

    \item \textbf{Solution 20:} Formula $A = \frac{1}{2} \sum (x_i y_{i+1} - x_{i+1}y_i)$.
    $A = \frac{1}{2} [ (0(0)-2(0)) + (2(3)-3(0)) + (3(4)-1(3)) + (1(0)-0(4)) ] = \frac{1}{2}[0 + 6 + 9 + 0] = \mathbf{7.5}$.

    \item \textbf{Solution 21:} $A = \pi/2$. Bottom boundary $y=0, dy=0 \implies 0$.
    Upper arc $x = \cos t, y = \sin t, dy = \cos t \,dt$.
    $\bar{x} = \frac{1}{\pi} \int_0^\pi \cos^2 t (\cos t) \,dt = \frac{1}{\pi} \int_0^\pi \cos^3 t \,dt = \mathbf{0}$. (Matches symmetry).

    \item \textbf{Solution 22:} Parameterize ellipse: $x = 3 \cos t, y = 4 \sin t$.
    $A = \frac{1}{2} \int_0^{2\pi} (x y' - y x') \,dt = \frac{1}{2} \int_0^{2\pi} (3 \cos t(4 \cos t) - 4 \sin t(-3 \sin t)) \,dt = \frac{1}{2} \int_0^{2\pi} 12 \,dt = \mathbf{12\pi}$.

    \item \textbf{Solution 23:} For $[0,a] \times [0,b]$, only the top boundary $(a,b) \to (0,b)$ contributes to $\oint y^3 \,dx$ (since $dy$ lines have $dx=0$, and $y=0$ line is 0).
    $I_x = -\frac{1}{3} \int_a^0 b^3 \,dx = -\frac{1}{3} [b^3 x]_a^0 = -\frac{1}{3}(0 - ab^3) = \mathbf{\frac{ab^3}{3}}$.

    \item \textbf{Solution 24:} \textit{Singularity Trap.} Field is undefined at $(0,0)$. 
    Create a unit circle $C_r$ centered at origin. By extended Green's theorem, $\oint_C = \oint_{C_r}$. Over $C_r$, $x = \cos t, y = \sin t$. $\int_0^{2\pi} \frac{-\sin t(-\sin t dt) + \cos t(\cos t dt)}{1} = \int_0^{2\pi} (\sin^2 t + \cos^2 t) \,dt = \mathbf{2\pi}$.

    \item \textbf{Solution 25:} Split integral: $\oint \frac{-y}{x^2+y^2}dx + \frac{x}{x^2+y^2}dy \quad + \quad \oint 2ydx + 3xdy$.
    First part $= 2\pi$ (singularity inside). Second part: $Q_x - P_y = 3 - 2 = 1 \implies \iint 1 dA = 9\pi$. Total $= \mathbf{11\pi}$.

    \item \textbf{Solution 26:} Singularity is at $(1,0)$, which is inside $x^2+y^2=16$.
    Since this is just a shifted fundamental circulation field, the line integral evaluates purely to $\mathbf{2\pi}$.

    \item \textbf{Solution 27:} $Q_x - P_y = 2xy^2 + 1 - (-1) = 2xy^2 + 2$.
    $\int_0^2 \int_{-1}^1 (2xy^2 + 2) \,dx \,dy$. The function $2xy^2$ is an odd function of $x$ over the symmetric interval $[-1,1]$, so it integrates to $0$. We are left with $\int_0^2 \int_{-1}^1 2 \,dx \,dy = \mathbf{8}$.

    \item \textbf{Solution 28:} $Q_x - P_y = 3x^2y - 2y$. 
    The region is symmetric across the x-axis ($y$ goes from $-2$ to $2$). The integrand is an odd function with respect to $y$ ($f(x,-y) = -f(x,y)$), and the boundaries depend only on even powers of $y$. Hence, the double integral evaluates identically to $\mathbf{0}$.

    \item \textbf{Solution 29:} Astroid area: $x = a \cos^3 t, y = a \sin^3 t$.
    $dx = -3a \cos^2 t \sin t \,dt, dy = 3a \sin^2 t \cos t \,dt$.
    $x dy - y dx = 3a^2 \sin^2 t \cos^2 t (\sin^2 t + \cos^2 t) = 3a^2 \sin^2 t \cos^2 t$.
    $A = \frac{1}{2} \int_0^{2\pi} 3a^2 (\sin t \cos t)^2 \,dt = \frac{3a^2}{8} \int_0^{2\pi} \sin^2(2t) \,dt = \frac{3a^2}{16} \int_0^{2\pi} (1 - \cos 4t) \,dt = \mathbf{\frac{3\pi a^2}{8}}$.

    \item \textbf{Solution 30:} Cycloid area using boundary formula $A = - \oint y \,dx$. 
    On the x-axis, $y=0 \implies 0$. For the arch (traced right to left to remain CCW, so $t$ from $2\pi$ to $0$):
    $A = - \int_{2\pi}^0 R(1-\cos t) \cdot R(1-\cos t) \,dt = R^2 \int_0^{2\pi} (1 - 2\cos t + \cos^2 t) \,dt = R^2(2\pi - 0 + \pi) = \mathbf{3\pi R^2}$.
\end{itemize}

\newpage
\section{Part 3: Concept Checklist \& Question Mapping}

Use this checklist to ensure you have reviewed all the critical concepts, traps, and shortcuts required for mastery of Green's Theorem.

\begin{itemize}[label=$\square$, leftmargin=*]
    \item \textbf{C1: Direct Evaluation vs. Theorem Verification} 
    \textit{Concept:} Ability to manually parameterize boundaries and sum segments to confirm the theorem's equivalence.
    \textit{Covered in Questions:} 1, 2

    \item \textbf{C2: Standard Cartesian Setup}
    \textit{Concept:} Accurately computing $Q_x - P_y$ and setting up standard Type I / Type II double integrals bounded by lines or parabolas.
    \textit{Covered in Questions:} 3, 4, 5, 6, 7

    \item \textbf{C3: Polar Transformation Setup}
    \textit{Concept:} Recognizing circular, semi-circular, or annular regions, swapping to $r$ and $\theta$, and never forgetting the Jacobian $r$.
    \textit{Covered in Questions:} 8, 9, 10, 11, 12

    \item \textbf{C4: Clockwise / Negative Orientation}
    \textit{Concept:} Visualizing the path trace. If it behaves like a clock, manually appending a MINUS SIGN to the Green's Theorem double integral.
    \textit{Covered in Questions:} 13, 14, 15, 16

    \item \textbf{C5: The "Area Shortcut"}
    \textit{Concept:} Halting complex integration if $Q_x - P_y = K$ (a constant). Evaluating simply as $K \times (\text{Area of Region})$.
    \textit{Covered in Questions:} 17, 18, 19

    \item \textbf{C6: Geometric Properties (Area, Centroids, MOI)}
    \textit{Concept:} Utilizing specialized line integral formulas (like $1/2 \oint xdy - ydx$) to uncover geometric properties purely from the boundary trace.
    \textit{Covered in Questions:} 20, 21, 22, 23

    \item \textbf{C7: Extended Green's Theorem \& Singularities}
    \textit{Concept:} Recognizing when $P$ or $Q$ are undefined inside the region (e.g., division by zero at origin). Bypassing via inner-circle parameterization.
    \textit{Covered in Questions:} 24, 25, 26

    \item \textbf{C8: Exploiting Odd Functions \& Symmetries}
    \textit{Concept:} Using odd function logic over symmetric intervals (like $[-a, a]$) to immediately evaluate nasty algebraic integrals to zero.
    \textit{Covered in Questions:} 27, 28

    \item \textbf{C9: Advanced Parameterization (Untested Shapes)}
    \textit{Concept:} Finding areas of complex non-trivial shapes (like Astroids or Cycloids) relying strictly on pure Calculus II trigonometric parameterization steps.
    \textit{Covered in Questions:} 29, 30
\end{itemize}

\end{document}